\chapter{The Epistemology of Reflexive Constraint Realism}
\label{ch:epistemology}

\begin{chapterguide}
This chapter gives \RCR\ an epistemology. It defines evidence, states
inferential rules, explains how belief should change, and proposes criteria for
comparing theories. The framework does not replace Bayesian, causal,
experimental, or interpretive methods; it specifies how their results can be
assembled into an objective but fallible claim.
\end{chapterguide}

The ontological claim that reality constrains inquiry is not enough. A world can
be real while our beliefs about it are systematically wrong. \RCR's
epistemology therefore asks how constraints become reasons. A reason is not
simply an event that occurs before belief. It is evidence when the event
discriminates among claims under a defensible measurement and inferential
model.

\section{Claims, targets, and alternatives}

Every scientific claim should be represented with five fields:

\begin{enumerate}
\item \textbf{Target:} the entity, process, relation, or population at issue;
\item \textbf{content:} what is being asserted about the target;
\item \textbf{scope:} the time, place, scale, population, and intervention
conditions;
\item \textbf{alternatives:} rival explanations or models;
\item \textbf{warrant:} the evidence and reasoning supporting the claim.
\end{enumerate}

A sentence such as \enquote{\(X\) causes \(Y\)} is incomplete. It does not say
which \(X\), which \(Y\), in what population, by what intervention, or against
which confounders. A scoped claim can be tested. A vague claim can survive by
changing its meaning.

\section{Evidence as constraint}

Evidence \(E\) supports a claim \(H\) when \(E\) is more expected under \(H\)
than under relevant alternatives, or when \(E\) forces a revision of those
alternatives. The Bayesian expression is the likelihood ratio:

\[
  \operatorname{LR}(E;H_1,H_0)
  =
  \frac{\Prob(E\mid H_1)}{\Prob(E\mid H_0)}.
\]

A likelihood ratio greater than one favours \(H_1\) over \(H_0\); a ratio less
than one favours \(H_0\). The ratio does not establish either hypothesis as
true. It measures the evidential direction and strength relative to a model.

In practice, \(\Prob(E\mid H)\) includes assumptions about sampling,
measurement, missingness, and analysis. \RCR\ therefore requires that the
evidence statement carry its measurement path and scope conditions with it.

\begin{principlebox}
\textbf{Evidence rule.} Do not ask whether a result is evidence in isolation.
Ask: evidence for which claim, against which alternative, under which
measurement path, and with what expected error if the claim were false?
\end{principlebox}

\section{The inferential rules}

The following rules are the operational core of \RCR. They are norms for
reasoning, not algorithms that mechanically settle every case.

\begin{enumerate}[label=\textbf{R\arabic*.}]
\item \textbf{Scope before strength.} State the claim's target and scope before
assigning confidence. A highly precise claim with an undefined target has
little epistemic content.
\item \textbf{Path traceability.} Trace the path from target to datum to
conclusion. Separate observation, transformation, model assumption, and
interpretation.
\item \textbf{Alternative exposure.} Compare a claim with serious alternatives.
A result that fits \(H\) but also fits many rivals has limited discriminating
force.
\item \textbf{Risk preference.} Prefer tests that create outcomes likely to
differ across hypotheses. Confirmation is stronger when the test was not
designed to guarantee success.
\item \textbf{Invariance search.} Vary observers, instruments, operational
definitions, samples, models, or interventions where the claim says the
result should remain stable.
\item \textbf{Causal separation.} Do not infer causal adequacy from predictive
fit alone. Require design, intervention, mechanism, or a clearly defended
counterfactual model.
\item \textbf{Transport discipline.} Do not move a result to a new setting
without identifying shared mechanisms and checking changed distributions.
\item \textbf{Commitment proportionality.} Match ontological commitment to the
strongest layer of evidence actually established.
\item \textbf{Revision readiness.} State what evidence would lower confidence or
change the model. A claim insulated from revision is not strengthened by
confidence.
\item \textbf{Social accountability.} Treat competent, relevant criticism as a
source of information about hidden assumptions, not merely as a political
obstacle.
\end{enumerate}

The rules can conflict. A high-risk test may be unethical. An exact replication
may be impossible. A new operationalization may alter the target. In such
cases, the conflict should be reported rather than hidden. \RCR\ is a framework
for making trade-offs visible, not a claim that every trade-off has one
correct answer.

\section{Belief revision}

Belief revision should be continuous, comparative, and calibrated. Continuous
does not mean that all beliefs must be represented by exact numbers. It means
that the language of confidence should reflect degrees of support. Comparative
means that confidence in one claim depends on how it performs against
alternatives. Calibrated means that predictions are checked against outcomes.

A useful reporting format separates:

\begin{itemize}
\item prior background assumptions;
\item new evidence;
\item the likelihood or expected error under each alternative;
\item the updated confidence;
\item remaining uncertainty and scope limitations.
\end{itemize}

This format is compatible with Bayesian updating, confidence intervals,
likelihoodist reasoning, error statistics, or qualitative comparative
inference. The common requirement is that the route from evidence to belief be
auditable.

\section{Theory comparison}

Theories should not be ranked by one virtue alone. A simple theory can be too
simple. A detailed theory can overfit. A highly predictive theory can lack a
causal mechanism. A metaphysically rich theory can make no risky difference.

\RCR\ proposes a comparative profile with seven dimensions:

\begin{enumerate}
\item empirical fit to established phenomena;
\item novel or risky success;
\item structural and causal coherence;
\item robustness under independent changes;
\item explanatory reach and integration;
\item simplicity relative to scope and complexity;
\item transparency, corrigibility, and transportability.
\end{enumerate}

For an informal representation, let a theory \(T\) have profile
\[
  \mathbf{V}(T) =
  (F,N,C,R,X,S,O),
\]
where the letters stand for fit, novelty, coherence, robustness, explanatory
reach, simplicity, and objectivity of process. The profile is not a single
score. In some decisions one dimension can dominate: a medical treatment may
require safety even if its explanatory theory is incomplete. In others,
structural coherence or novel prediction may matter more.

A weighted score can be used when a purpose requires one:
\[
  Q(T) = \sum_{j=1}^{7} w_j V_j(T),
  \qquad
  w_j \geq 0,\quad \sum_{j=1}^{7}w_j=1.
\]
Here \(w_j\) are explicit weights chosen for a stated purpose. The equation
does not create objectivity by arithmetic. It prevents hidden trade-offs from
masquerading as neutral judgment.

\begin{table}[ht]
\centering
\smalltable
\begin{tabularx}{\textwidth}{p{0.22\textwidth}Y Y}
\toprule
Dimension & Strong performance & Typical warning sign\\
\midrule
Fit & Accounts for reliable phenomena within error & Selective use of data\\
Novelty & Predicts or reveals outcomes not used to tune the theory & Post hoc accommodation\\
Coherence & Links mechanisms, equations, and concepts without contradiction & Ad hoc patches\\
Robustness & Survives instruments, samples, and model changes & Fragility to one pipeline\\
Reach & Connects domains or scales with clear bridges & Vague unification\\
Simplicity & Achieves scope without needless complexity & Underfitting or oversimplification\\
Process objectivity & Makes criticism and revision possible & Opaque methods or status protection\\
\bottomrule
\end{tabularx}
\caption{A comparative profile for theory choice. No dimension is sufficient on
its own.}
\label{tab:theory-profile}
\end{table}

\section{Explanation and evidence are distinct}

A theory may explain a phenomenon but fit data poorly because the measurements
are noisy. A statistical model may fit and predict accurately without
identifying the mechanism. A mechanism may be true while a quantitative
parameter is wrong. \RCR\ avoids forcing these into one score by reporting
separate warrants.

This distinction is also useful for scientific communication. A paper should
not describe an associational analysis as a causal explanation merely because
the proposed story is plausible. It should not describe a computational
possibility as an empirical prediction until the bridge to the target is
identified. It should not describe a normatively preferred outcome as a
scientific conclusion without separating values from evidence.

\section{Objectivity criteria}

A claim qualifies as objective in the \RCR\ sense when it satisfies the
following criteria to a degree appropriate to its field:

\begin{enumerate}
\item \textbf{World resistance:} the claim confronts effects that investigators
cannot freely choose.
\item \textbf{Cross-perspective convergence:} relevantly different observers or
methods identify the same pattern.
\item \textbf{Measurement accountability:} the path from target to result is
traceable and error-aware.
\item \textbf{Intervention or counterfactual grip:} the claim supports a
well-defined change relation where one is possible.
\item \textbf{Model transparency:} idealizations and bridge assumptions are
stated.
\item \textbf{Critical independence:} competent others can inspect and challenge
the claim.
\item \textbf{Scope honesty:} the claim does not travel further than its
evidence.
\item \textbf{Revision capacity:} the inquiry can detect and incorporate error.
\end{enumerate}

These criteria are not a checklist that produces certainty. They are dimensions
along which a research system can improve. A claim may score strongly on
measurement and weakly on intervention because the target is historical. It may
score strongly on causal grip and weakly on transport because the setting is
unique.

\section{What makes the epistemology reflexive}

Reflexivity means more than adding a disclaimer about bias. It requires
studying how the inquiry's own categories and procedures change the target,
the evidence, or the community. In social science, a classification can
become an incentive. In medicine, a diagnostic label can affect treatment. In
physics, a detector design defines which events are accessible. In all cases,
the inquiry should ask what it has made visible and what it has excluded.

A reflexive epistemology is not self-defeating. The fact that an inquiry has
conditions does not imply that its conclusions are only about those conditions.
It means that a conclusion earns broader authority when it survives changes to
the conditions that are not supposed to matter.

\section*{Chapter summary}

\RCR\ treats evidence as a claim-relative constraint, not as a free-floating
fact. Its ten inferential rules require scope, traceability, alternatives,
risk, invariance, causal separation, transport discipline, proportional
commitment, revision readiness, and social accountability. Theory comparison
uses a profile of fit, novelty, coherence, robustness, reach, simplicity, and
process objectivity. The criteria of objectivity are graded, field-sensitive,
and designed to make error discoverable.

\begin{discussion}
\begin{enumerate}
\item Should a theory with stronger fit but weaker causal explanation be
preferred for prediction?
\item Which of the ten rules is most difficult to satisfy in social science?
\item Can explicit values improve theory comparison rather than bias it?
\end{enumerate}
\end{discussion}

\begin{furtherreading}
Compare the updating logic here with Bayes (1763), Reichenbach (1938), Mayo
(1996), and the causal frameworks of Pearl (2009) and Woodward (2003).
\end{furtherreading}

